% ===================================================================
%  اللوحة رقم ١٦ — المحلّ الكبير. --- البازار. اللعب.
%
%  Traduction, faite en 2026, en arabe levantin d'aujourd'hui, sur
%  l'ido de texto/io. Regles de la colonne : voir 10-lawha-01.tex.
%  Une seule numerotation. 33 blocs, 106 renvois. Dernier tableau du
%  livret. UNE note d'auteur, t16-noto-1, posee entre les blocs 8 et
%  9 ; son appel se compose en parentheses simples, « (*) », dans le
%  bloc 13. L'ido coupe cinq alineas sur des « suite » de feuillet
%  (3, 5, 8, 13, 16) ; la traduction n'herite d'aucune des cinq et
%  ecrit un bloc chacun.
%
%  ECART DECLARE : t16-15-1. L'ido saute la lettre « (d) » que le
%  fac-simile grave pourtant ; renvoji.py laisse donc ce bloc passer
%  SANS RIEN VERIFIER. L'ordre a ete compte a la main : 85, a, b, c,
%  d, e, f, 86, 87, 88, 89 — onze appels.
%
%  LE TITRE N'A PAS BESOIN D'ETRE TRADUIT, ET C'EST LA MEME
%  DECOUVERTE QUE LA COLONNE GUJARATIE AU MEME TABLEAU. « La Bazaro »
%  s'ecrit ici « البازار », et la colonne gujaratie ecrivait
%  « બજાર » — un seul mot persan, bāzār, parti de Chiraz vers l'est
%  et vers l'ouest, arrive en Inde par les Moghols, au Levant par les
%  Seldjoukides, en France par l'italien. Le francais de Rochelle
%  ecrit « bazar » et ne s'apercoit pas qu'il ecrit du persan. TROIS
%  COLONNES ET L'ORIGINAL SUR UN SEUL MOT, ET C'EST LE DERNIER TITRE
%  DU LIVRET.
%
%  L'ALINEA 5 EST LE SEUL ENDROIT DE CETTE COLONNE OU LE LEVANTIN
%  N'A PAS DE REGISTRE A LUI, ET IL FAUT L'ECRIRE PLUTOT QUE DE FAIRE
%  SEMBLANT. Quinze tableaux durant, la regle de la colonne a ete de
%  ne pas remonter vers l'arabe de l'ecole. La bataille de soldats de
%  plomb du grand-pere ne le permet pas : le siege, le bombardement,
%  l'assaut, la sortie, la retraite, l'armistice, le traite de paix
%  n'ont AUCUNE forme dialectale. On ne raconte pas une guerre en
%  levantin — on la raconte en fusha, parce que c'est la langue des
%  journaux et des communiques, et un Damascene qui parlerait de
%  siege a table dirait حصار comme le presentateur. La colonne ecrit
%  donc l'alinea 5 en standard, exprès, et le dit ici : UNE REGLE
%  QU'ON SUSPEND EN LE DISANT VAUT MIEUX QU'UNE REGLE QU'ON TRAHIT EN
%  SILENCE.
%
%  LE CROCODILE (12) EST « التمساح », ET C'EST LA TROISIEME FOIS QUE
%  CETTE COLONNE DESCEND SOUS L'ARABE. Le mot vient du copte ⲉⲙⲥⲁϩ,
%  et le copte de l'egyptien ancien : il est entre en arabe avec la
%  conquete et n'en est jamais ressorti. Le releve tenait deja تموز
%  au tableau 5, akkadien, et صمّون au tableau 15, akkadien aussi par
%  le grec et l'arameen. Trois substrats, trois tableaux : le temps,
%  la nourriture, et pour finir une bete du Nil. Deux renvois plus
%  loin, l'hippopotame (22) est « فرس النهر », qui est le calque
%  mot-pour-mot du grec ἱπποπόταμος fait par les traducteurs de
%  Bagdad — celui-la n'est pas un substrat mais une traduction, et il
%  a mille deux cents ans.
%
%  ET LA DERNIERE TENTATION DE LA COLONNE A ETE REFUSEE. Le theatre
%  de marionnettes (41) et ses trois personnages — Arlequin (42),
%  Polichinelle (43), Pierrot (44) — appellent de toutes leurs forces
%  « كراكوز », le Karagöz du theatre d'ombres ottoman, qui se jouait
%  dans les cafes de Damas et d'Alep jusqu'au siecle dernier et que
%  tout le monde connait. Ce n'est pas ce que la planche montre. On
%  ecrit donc « مسرح العرايس » et les trois noms transcrits. ON
%  MODERNISE LA LANGUE, JAMAIS LES CHOSES — et la regle se verifie
%  mieux au dernier tableau qu'au premier, parce qu'au dernier on
%  sait exactement ce qu'on perd.
%
%  Enfin l'eclipse (b) demande deux mots la ou le francais n'en a
%  qu'un : خسوف pour la lune, كسوف pour le soleil. La langue ne les
%  confond jamais, et le renvoi unique de l'ido en couvre donc deux.
% ===================================================================

%%K t16-apar-1 apar
\VUsaut{4.0mm}
\VUtitre{15.0pt}{60}{اللوحة رقم ١٦}
\VUsaut{2.0mm}
\VUpk{13.2pt}{40}{المحلّ الكبير. --- البازار.}
\VUpk{13.2pt}{40}{اللعب.}
\VUsaut{3.4mm}

%%K t16-01-1 p
1. --- السنة الجديدة قرّبت. والمحلّات الكبيرة حضّرت بسطاتها
\VUgras{لعب} \textsuperscript{(1)} و\VUgras{هدايا} رأس السنة.

%%K t16-01-2 p
وأنا طلبت من جدّي ياخدني عالبازار لشوف هالمعرض الحلو، وهو وافق.

%%K t16-02-1 p
2. --- أوّل شي وقفنا قدّام أطقم سلاح حلوة، فيها \VUgras{بارودة}،
و\VUgras{قبّعة عسكريّة}، و\VUgras{سيف}، و\VUgras{حزام سيف} وجوز
\VUgras{كتفيّات}، أو \VUgras{خوذة} و\VUgras{درع} \textsuperscript{(3)}
و\VUgras{مسدّس} \textsuperscript{(4)}. وكلّ هاد شدّني أكتر بكتير من
\VUgras{الفوانيس السحريّة} \textsuperscript{(2)}.

%%K t16-tit-1 sub
\VUsaut{3.0mm}
\VUpk{10.2pt}{40}{مناظر حلوة.}
\VUsaut{2.6mm}

%%K t16-03-1 p
3. --- وبنفس القسم كان في \VUgras{خفير} \textsuperscript{(5)}
بـ\VUgras{كشك حراسته} ؛ مبيّن عليه عم يحرس قدّام حديقة حيوانات كاملة
من كرتون. وشو ما كان في حيوانات هونيك : \VUgras{ببغا}
\textsuperscript{(6)} على عصايته، و\VUgras{وحيد قرن} \textsuperscript{(7)}
وقرنه على منخاره، و\VUgras{رنّة} \textsuperscript{(8)}، و\VUgras{نعامة}
\textsuperscript{(9)}، و\VUgras{قرد} \textsuperscript{(10)} عم يكسّر جوزة،
و\VUgras{تعلب} \textsuperscript{(11)} شايل جاجة، و\VUgras{تمساح}
\textsuperscript{(12)} فاتح فكّيه الكبار، و\VUgras{جمل}
\textsuperscript{(13)}، و\VUgras{سلوقي} \textsuperscript{(14)} أنيق،
و\VUgras{فهد} \textsuperscript{(15)}، وبعدين حيوانين ما كنت بعرفهن وطلعوا،
على ما قالولي، \VUgras{ابن عرس} \textsuperscript{(16)} و\VUgras{ضبع}
\textsuperscript{(17)}. وكمان كان في \VUgras{نمر} \textsuperscript{(18)}
و\VUgras{ديب} \textsuperscript{(21)} ؛ وكتير جنبهن كان في \VUgras{جرذان}
\textsuperscript{(19)} زغار حلوين و\VUgras{فيران} \textsuperscript{(20)}
زغيرة كتير من مطّاط. وبالآخر \VUgras{فرس نهر} \textsuperscript{(22)}
و\VUgras{جمل وحيد السنام} \textsuperscript{(23)} أحدب سكّروا المجموعة.
وكنت بضلّ هونيك ساعات كمان لو ما كان جنبي مخيّم رائع معبّى
\VUgras{عساكر رصاص} \textsuperscript{(29)} شدّ انتباهي.

%%K t16-tit-2 sub
\VUsaut{4.0mm}
\VUpk{10.2pt}{40}{عساكر وحرب.}
\VUsaut{2.8mm}

%%K t16-04-1 p
4. --- وجدّي، اللي هو \VUgras{معاق حرب} \textsuperscript{(24)} واللي
اضطرّ يترك الجيش بسبب \VUgras{جرح} جابله \VUgras{الميداليّة
العسكريّة}، بيحبّ كتير يحكيلي عن \VUgras{الحملات} اللي شارك فيها. وكان
مبسوط وهو عم يشرحلي حركات الجيش الصغير المصغّر. وشلّة عساكر حقيقيّين
كانوا عم يتفرّجوا عالبازار : \VUgras{جندي هندسة} \textsuperscript{(25)}،
و\VUgras{صيّاد جبلي} \textsuperscript{(26)} لابس \VUgras{بيريه}،
و\VUgras{رقيب أوّل} \textsuperscript{(27)} من المشاة و\VUgras{رقيب
مدفعيّة} \textsuperscript{(28)} و\VUgras{شريط} دهب على كمّه، وقفوا جنبنا
ليسمعوا الشرح.

%%K t16-05-1 p
5. --- وجدّي، وهو مفتخر كتير لأنّه صار عنده سامعين متل هدول، ارتجل
خطّة معركة كاملة.

%%K t16-05-2 p
المدينة المحصّنة، بـ\VUgras{أسوارها} و\VUgras{خنادقها}، كان
\VUgras{الجيش المعادي} محاصرها، وبعد \VUgras{قصف} طويل صار جاهز
لـ\VUgras{الاقتحام}. بس \VUgras{المحاصَرين}، واللي الجوع ضيّقهم، جرّبوا
يعملوا \VUgras{اقتحام معاكس} على \VUgras{مخيّم} \VUgras{المحاصِرين}.
وبعد ما ردّوا \VUgras{الهجوم} بـ\VUgras{قتال} عنيد حوالين
\VUgras{الخيم}، المحاصِرين \VUgras{المنتصرين} رموا \VUgras{أسرابهم}
لـ\VUgras{ملاحقة} المهاجمين \VUgras{المهزومين}. وهدول
\VUgras{انسحبوا}، وغطّوا \VUgras{ساحة المعركة} بـ\VUgras{قتلى}
و\VUgras{جرحى}. و\VUgras{حمّالو النقّالات} شالوا الجرحى
عالـ\VUgras{إسعاف} ليعالجهن \VUgras{الجرّاح}.

%%K t16-05-3 p
ورغم المقاومة البطوليّة تبع كم \VUgras{كتيبة} تشكّلوا بـ\VUgras{مربّع}،
كانت \VUgras{الهزيمة} كاملة، وباقي الجيش \VUgras{هرب}، وترك
\VUgras{أسرى} كتار. والعدوّ رايح يتوّج \VUgras{انتصاره} باحتلال المدينة.
يمكن هاي تكون نهاية الحملة، وبعد \VUgras{هدنة}، بيصير ممكن يوقّعوا
\VUgras{معاهدة سلام}.

%%K t16-tit-3 sub
\VUsaut{4.4mm}
\VUpk{10.2pt}{40}{هدايا رأس السنة.}
\VUsaut{2.8mm}

%%K t16-06-1 p
6. --- والعساكر الشباب، اللي ضحكوا وهنّي عم يسمعوا حكاية المحارب
القديم، طلبوا منّه كم شرح تاني. أمّا أنا فلفّيت بالمحلّ لتفرّج على
\VUgras{نشّابة} \textsuperscript{(32)} حلوة و\VUgras{قوس}
\textsuperscript{(33)} مع \VUgras{سهمه} \textsuperscript{(31)}. ولقيت
هونيك مجموعة حيوانات تانية : \VUgras{طيرين مغرّدين}
\textsuperscript{(34)} : \VUgras{بلبل} و\VUgras{هازجة} أوتوماتيك عم
يغنّوا بملء حلقهن، وسلسلة حيوانات غريبة : \VUgras{وطواط}
\textsuperscript{(35)}، و\VUgras{أفعى بوا} \textsuperscript{(36)}،
و\VUgras{سلحفاة} \textsuperscript{(37)}، و\VUgras{فقمة}
\textsuperscript{(38)}، و\VUgras{حوت} \textsuperscript{(39)}
و\VUgras{قرش} \textsuperscript{(40)} من مطّاط منفوخ.

%%K t16-07-1 p
7. --- وما وقفت كتير لتفرّج عليهن، لأنّي شفت اللي كنت بحلم فيه :
\VUgras{مسرح عرايس} \textsuperscript{(41)} رائع مع \VUgras{ديكورات}
فخمة و\VUgras{ستارة} حمرا بتاخد العقل. وعلى \VUgras{الخشبة}، تنين
\VUgras{ممثّلين} مبيّن عليهن عم يلعبوا \VUgras{دور} بـ\VUgras{مسرحيّة}
كتير حلوة، لأنّ \VUgras{المتفرّجين} الزغار من كرتون، اللي محطوطين
بالـ\VUgras{مقصورات} أو قاعدين عالـ\VUgras{كراسي} وبـ\VUgras{الصالة}
ورا \VUgras{قائد الأوركسترا}، كانوا عم يصفّقوا بحماس.

%%K t16-07-2 p
بس للأسف، هالمسرح كان غالي كتير.

%%K t16-08-1 p
8. --- وبعدين، اختيار اللعب كان صعب، لأنّ بالفاترينات كان متكوّم كلّ
نوع لعب : \VUgras{أرلوكان} \textsuperscript{(42)}، و\VUgras{بولشينيل}
\textsuperscript{(43)}، و\VUgras{بييرو} \textsuperscript{(44)}،
و\VUgras{بوم} \textsuperscript{(45)} عم يرفرف بجناحاته، و\VUgras{شحارير}
\textsuperscript{(46)} بتصفّر، و\VUgras{كلاب بودل} بتنبح،
و\VUgras{فونوغراف} \textsuperscript{(47)} و\VUgras{ألعاب صبر.} وفي وحدة
من هاللعب سلّتني كتير : كانت بتمثّل \VUgras{معركة بحريّة}
\textsuperscript{(48)}، فيها \VUgras{مركب} عم يهجموا عليه
\VUgras{قراصنة} جنب برّ مزروع \VUgras{نخل جوز الهند.}

%%K t16-noto-1 noto
\VUnotes{143.2mm}{%
\textit{(*) شوف اللوحة رقم ٦.}%
}

%%K t16-09-1 p
9. --- وقدرت أتفرّج بكلّ راحة على هالعجايب كلّها، لأنّه ما كان في
غير شوي ناس بالمحلّ، بالكاد كم واحد عم يتمشّى، و\VUgras{جندي دراغون}
\textsuperscript{(49)} و\VUgras{محصّل} \textsuperscript{(50)} كان عم يقدّم
\VUgras{كمبيالة} عند \VUgras{الصندوق} \textsuperscript{(51)} الرئيسي.

%%K t16-10-1 p
10. --- وسألت جدّي شو فايدة الكتب المركونة عالرفوف ورا
\VUgras{الصرّافة} ؛ جاوبني إنّهن \VUgras{دفاتر الحسابات}.

%%K t16-tit-4 sub
\VUsaut{3.6mm}
\VUpk{10.2pt}{40}{تفرّج عالدكاكين.}
\VUsaut{2.8mm}

%%K t16-11-1 p
11. --- ولمّا تركنا قسم اللعب، رحنا على قسم السروج لنشوف
\VUgras{السروج} \textsuperscript{(53)}، و\VUgras{الركابات}
\textsuperscript{(54)}، و\VUgras{اللجم} \textsuperscript{(55)}،
و\VUgras{الشكايم} \textsuperscript{(56)} و\VUgras{الكرابيج}. ووقت كان
\VUgras{رئيس القسم} \textsuperscript{(57)} عم يعطي جدّي كم معلومة، رحت
أتفرّج على بسطة \VUgras{الصيني} و\VUgras{الكريستال}
\textsuperscript{(58)} اللي فيها \VUgras{سلطانيّات} حلوة
و\VUgras{أباريق شاي} \textsuperscript{(59)} رفيعة.

%%K t16-12-1 p
12. --- ولنطلع عالطابق الأوّل، مرقنا ورا فاترينة
\VUgras{الكورسيهات} \textsuperscript{(60)}، جنب قسم الجوارب، اللي كان
معروض فيه \VUgras{لباسات} \textsuperscript{(63)} كتير حلوة. وجنبها،
\VUgras{بيّاعة} \textsuperscript{(61)} كانت عم تخلّي \VUgras{شارية}
\textsuperscript{(62)} تنقّي \VUgras{أقمشة} \textsuperscript{(64)} حرير
حلوة.

%%K t16-12-2 p
ووقت كنّا عم نطلع عالدرج، لاحظت إنّ المحلّ مزيّن
بـ\VUgras{فوانيس} \textsuperscript{(65)} يابانيّة، و\VUgras{شماسي}
\textsuperscript{(66)}، و\VUgras{نخل} \textsuperscript{(70)} كبير كتير.

%%K t16-13-1 p
13. --- وعالطابق الأوّل، ما كان في شي كتير لنشوفه ؛ بس مفروشات (*)،
و\VUgras{خزنات حديد} \textsuperscript{(67)}، و\VUgras{خزانات جوارير}
\textsuperscript{(68)} و\VUgras{عربايات أطفال} \textsuperscript{(69)}. بس
منظر المحلّ كلّه كان \VUgras{كتير مسلّي}.

%%K t16-14-1 p
14. --- وتحت منّا، شفت آلات الموسيقى : \VUgras{قيثارات}
\textsuperscript{(71)}، و\VUgras{غيتارات} \textsuperscript{(72)}،
و\VUgras{مندولينات} \textsuperscript{(73)}، و\VUgras{أبواق بصمّامات}
\textsuperscript{(74)}، و\VUgras{كلارينيتات} \textsuperscript{(75)}،
وكمنجات مع \VUgras{قوسها} \textsuperscript{(76)} وحتّى \VUgras{طبل كبير}
\textsuperscript{(77)}. وأبعد شوي، عند قسم القرطاسيّة، \VUgras{طالب}
\textsuperscript{(78)} كان عم يفاصل \VUgras{البيّاع} \textsuperscript{(79)}
على شنتة دفاتر. وجنبهن، ميّزت \VUgras{كتاب ألفباء} \textsuperscript{(80)}
حلو.

%%K t16-14-2 p
وبنصّ المحلّ كانت منصوبة \VUgras{شجرة ميلاد} \textsuperscript{(81)}
عالية معبّاية شمعات وحاجات زغيرة وكلّ نوع لعب : \VUgras{خيل خشب}
\textsuperscript{(82)}، و\VUgras{عرايس} \textsuperscript{(83)}، إلخ.

%%K t16-14-3 p
و\VUgras{قسم العطورات} \textsuperscript{(84)} بـ\VUgras{غسولات
أسنانه} و\VUgras{عطوراته} المتنوّعة كان عم ينشر ريحة كتير طيّبة وصلت
لعنّا.

%%K t16-tit-5 sub
\VUsaut{3.4mm}
\VUpk{10.2pt}{40}{منشتري شي.}
\VUsaut{2.8mm}

%%K t16-15-1 p
15. --- ولأنّ جدّي بلّش يحسّ إنّ الوقت طوّل، نزلنا عالدرج الكبير.
ووقف شوي قدّام \VUgras{لوحة فلك} \textsuperscript{(85)} بتمثّل، متل ما
قال لي، \VUgras{مذنّبات} \textsuperscript{(a)}، و\VUgras{خسوف القمر
وكسوف الشمس} \textsuperscript{(b)}، وبوصلة رياح مع \VUgras{الجهات
الأربع} \textsuperscript{(c)} \VUgras{(الشمال، والجنوب، والشرق والغرب)}،
وخريطة \VUgras{نصفي الكرة} \textsuperscript{(d)}، و\VUgras{الكواكب}
\textsuperscript{(e)} و\VUgras{المجموعة الشمسيّة} \textsuperscript{(f)}.
وأنا ما فهمت شي من هاد كلّه، وكان يهمّني أكتر بكتير بدلة \VUgras{صبي
التوصيل} \textsuperscript{(86)} المشرّطة وهو مارق، و\VUgras{مراوح اليد}
\textsuperscript{(87)} الحلوة اللي كانت جنبي، جنب \VUgras{الجزدانات}
\textsuperscript{(88)} و\VUgras{المحافظ} \textsuperscript{(89)}.

%%K t16-16-1 p
16. --- وكمان انبهرت عالبسطة بـ\VUgras{الريش} \textsuperscript{(90)}،
و\VUgras{بزيم الأحزمة} \textsuperscript{(91)}، و\VUgras{الورد
الاصطناعي} \textsuperscript{(94)} الحلو المرتّب بشكل لطيف بسلال.
و\VUgras{البنفسج}، و\VUgras{المنثور}، و\VUgras{السنبل}
\textsuperscript{(95)}، و\VUgras{إبرة الراعي} \textsuperscript{(96)}،
و\VUgras{الزنبق} \textsuperscript{(97)} و\VUgras{أبو خنجر}
\textsuperscript{(98)} كانوا مقلّدين كتير منيح.

%%K t16-tit-6 sub
\VUsaut{4.4mm}
\VUpk{10.2pt}{40}{هديّة جدّي.}
\VUsaut{2.8mm}

%%K t16-17-1 p
17. --- وطلبت من جدّي يشتريلي وحدة من \VUgras{فرش الأسنان}
\textsuperscript{(92)} اللي كانت بعلبة جنب \VUgras{دبابيس الشعر}
\textsuperscript{(93)}. ووافق، ورحت بنفسي ادفع مشترايي عند الصندوق،
اللي كانوا عم يحضّروا جنبه \VUgras{إرساليّة} \textsuperscript{(100)}
كبيرة لـ\VUgras{زبون} \textsuperscript{(99)}.

%%K t16-18-1 p
18. --- وفجأة صار شوي هرج بين الناس الواقفين جنب \VUgras{التحف
البرونزيّة} \textsuperscript{(101)} الفنّيّة. \VUgras{حرامي}
\textsuperscript{(102)} كان لسّا منمسك بالجرم المشهود من \VUgras{مفتّش}
\textsuperscript{(103)}، وهو عم يحاول يسرق واحد. واهتمّوا يجيبوا شرطي
لـ\VUgras{توقيفه}، ووقت كنّا عم نطلع بالزبط، كانوا عم ياخدوا الجاني
عالحبس.
